\section{Cómo se identificaron las cancelaciones}\label{sec:cancelaciones}
Al añadir los pesos de dos cruces de un rectángulo, una misma línea puede aparecer en ambos.
Su vacancia se cuenta una sola vez. Aplicar esta regla a los dos órdenes del producto da
\(\Delta_1,\Delta_2\) (sección~\ref{sec:contracciones}). El método de rectángulos, adaptado
de~\cite[\S4]{tc}, suma todos los centros de cada patrón en la malla. La extrapolación espacial
se estudia en la sección~\ref{sec:rectangulos}.

Para el permanente, Divoux y Wang dan el determinante principal, McCullagh desarrolla
correcciones formales~\cite{dw,mcc} y Li trata perturbaciones de norma de línea sublineal~\cite{li}.
Aquí se controla el cambio del resto al retirar las dos filas y las dos columnas de un par.
Al reinsertar sus dos casillas se conserva el estado de los otros dos vértices. Así la media
del par aporta \(J\), y la interacción de las dos diagonales entra una sola vez en \(-S\)
(apéndices~\ref{app:boundary} y~\ref{app:weights}).

La última cancelación exige conservar el signo. Al comparar dos cortes del mismo tablero se
obtiene la convergencia de la integral truncada. Un término lineal en su numerador produciría
una divergencia logarítmica; la prueba del apéndice~\ref{app:closing} permite descartarlo.
El catálogo del suplemento conserva los ensayos que condujeron a estas tres identificaciones.
